% moderncv
 
\documentclass[11pt,a4paper]{moderncv}
\moderncvtheme[blue]{banking}
\usepackage[utf8]{inputenc}
\usepackage[top=1cm, bottom=1.5cm, left=1.5cm, right=1.5cm]{geometry}

%custom head & body
\usepackage{moderncvheadcustom}
\usepackage{moderncvbodycustom}

% Largeur de la colonne pour les dates
\setlength{\hintscolumnwidth}{2.5cm}

\firstname{Max}
\familyname{Mattone}
\title{Responsable de projets numériques / Product Owner}
\address{Chaumont}{France}
\email{mattone.m@gmail.com}
\mobile{+33 6 66 57 32 64}
\extrainfo{Permis B}

\begin{document}
\maketitle

\section{Formation}
\cventry{2011 -- 2014}{Diplôme d'Ingénieur}{Mines Douai -- École d'ingénieurs française}{Douai}{France}{Spécialisation en technologies de l'information et de la communication}
\cventry{2009 -- 2011}{L3 Physique et Ingénierie}{Lycée Masséna}{Nice}{France}{Classes Préparatoires aux Grandes Écoles (CPGE)}

\section{Expériences Professionnelles}

\cventry{juin 2025 -- aujourd'hui}{Consultant IT indépendant \& contributeur open source}{Indépendant}{}{}{
\begin{itemize}
    \item \textbf{Développement d'outils sur mesure :} Conception et développement d'un outil de réconciliation de facturation pour un client (bénévolement).
    \item \textbf{Open source (Pharo \& PharoJS) :} Contribution active à ces écosystèmes et maintien d'une implication technique continue.
\end{itemize}}

\cventry{nov 2020 -- juin 2025}{Ingénieur Solutions Client / Avant-vente}{BrowserStack}{Dublin}{Irlande}{
Ingénieur solutions client pour la plateforme de test BrowserStack sur le segment Enterprise EMEA.
\begin{itemize}
\setlength{\itemindent}{0.25in}
\item \textbf{Succès clients Enterprise :} Pilotage de la stratégie technique et de son exécution pour le compte de L'Oréal, avec une amélioration de 2x de l'efficacité des tests et une optimisation significative des coûts.
\item \textbf{Leadership proactif :} Identification et conduite autonome de cette initiative à forte valeur métier, au-delà du périmètre standard du poste, aboutissant à une étude de cas publiée.
\item Maintien d'une communication proactive et d'un alignement technique avec les clients Enterprise.
\item Accompagnement de l'intégration de la plateforme dans les cycles de développement clients.
\item Présentation des nouvelles fonctionnalités et produits, recueil des retours et stimulation de l'adoption.
\item Garantissement du succès technique lors des phases de Proof of Concept (PoC).
\end{itemize}
}

\cventry{juillet 2019 -- nov 2020}{Lead Engineer technique Cloud}{Orange, consultant Sogeti}{Sophia Antipolis}{France}{
Développement Backend pour le portail Orange.
\begin{itemize}
\setlength{\itemindent}{0.25in}
\item Conception, implémentation et maintenance de backends PHP.
\item Réconciliation de sources de données multiples via MySQL et Elasticsearch.
\item Intégration et déploiement continus (CI/CD) dans le cloud avec GitLab et CloudFoundry.
\item Monitoring de la production avec Kibana et Elasticsearch.
\end{itemize}
}

\cventry{sept 2017 -- juin 2019}{Ingénieur Consultant Logiciel}{Amadeus, consultant Sogeti}{Sophia Antipolis}{France}{
Mission de développement back-end au sein de l'équipe de distribution de réservations hôtelières.
\begin{itemize}
\setlength{\itemindent}{0.25in}
\item Équipe Scrum dans une organisation SAFe.
\item Développement backend C++ avec MoM (Message-Oriented Middleware).
\item Analyse d'incidents et maintenance des composants de la solution en intégration continue.
\item Études et implémentation de Proof of Concept (PoC) pour améliorer l'outillage interne.
\end{itemize}
}

\cventry{mai 2016 -- sept 2017}{Ingénieur Consultant QA Logiciel}{Air-France, consultant Sogeti}{Sophia Antipolis}{France}{
Mission d'ingénierie QA pour le projet Cascade --- programme de refonte du système de planification de vols d'Air-France.
\begin{itemize}
\setlength{\itemindent}{0.25in}
\item Analyse fonctionnelle des besoins clients : planification de vols standards, ajustements de diffusion, \ldots
\item Conception et exécution de suites de tests manuelles et automatiques en méthodologie Agile.
\item Garantie de la qualité de la suite de tests de non-régression.
\item Conception et développement d'un injecteur de messages interne (SOAP, SSM, ASM) pour les suites de tests automatisées.
\end{itemize}
}

\cventry{oct 2015 -- nov 2015}{Ingénieur de Recherche Logiciel}{École des Mines Douai}{Douai}{France}{
Étude et développement de stratégies pour l'exploration et la cartographie de couverture Wi-Fi.
Cette étude fait suite à un travail original de 2014 aux Mines Douai visant à relever le défi de la géolocalisation intérieure par la cartographie Wi-Fi domestique.
\begin{itemize}
\setlength{\itemindent}{0.25in}
    \item Conception de l'architecture pour l'intelligence robotique.
    \item Développement C++ pour la gestion des cartes de couverture.
    \item Développement Pharo (Smalltalk) des définitions de stratégies d'exploration.
    \item Conception de stratégies pour la cartographie Wi-Fi.
\end{itemize}
}

\cventry{nov 2014 -- oct 2015}{Ingénieur de Recherche Logiciel}{Inria Lille - Nord Europe \& École des Mines Douai}{Lille}{France}{
Comment tirer pleinement profit d'un environnement de live coding comme Pharo dans un système robotique hautement distribué ?
\begin{itemize}
\setlength{\itemindent}{0.25in}
    \item Développement d'un prototype pour diffuser des mises à jour dynamiques et des migrations au sein de la plateforme Pharo.
    \item Conception et implémentation d'un processus de Bootstrap pour la plateforme Pharo.
    \item Contribution au projet open-source PhaROS (client ROS (Robotic Operating System) pour Smalltalk).
    \item Contribution au projet open-source Pharo.
    \item Démonstration SLAM (Simultaneous Localization And Mapping) utilisant le Lidar et les robots robuLAB et Kobuki.
\end{itemize}
}

\subsection{Stages}

\cventry{avril -- oct 2014}{Ingénieur Logiciel}{Inria Lille - Nord Europe}{Lille}{France}{
Étude et implémentation d'un compilateur JIT (Just In Time) agnostique de l'architecture, en tant que stagiaire dans l'équipe RMoD.
\begin{itemize}
\setlength{\itemindent}{0.25in}
\item Étude et implémentation d'un compilateur multi-architecture (principalement x86 \& ARM).
\item Contribution au projet open-source Pharo.
\end{itemize}
}

\cventry{mars -- sept 2013}{Ingénieur Web}{Dynaclic}{Bruxelles}{Belgique}{
Développement d'une plateforme ERP pour écoles, collèges et universités.
\begin{itemize}
\setlength{\itemindent}{0.25in}
\item Conception et maintenance de bases de données MySQL.
\item Développement Web PHP avec le framework Yii.
\end{itemize}
}

\cventry{mai -- août 2012}{Assistant Ingénieur}{ESA (Agence spatiale européenne)}{Noordwijk}{Pays-Bas}{
Collaboration avec le segment de contrôle des vols habités sur le projet ERA (European Robotic Arm).
\begin{itemize}
\setlength{\itemindent}{0.25in}
\item Définition et documentation des interfaces réseau assurant des communications sécurisées pour opérer le bras robotique en toute sécurité.
\item Conception et implémentation d'une base de données pour faciliter la gestion de projet.
\end{itemize}
}

\section{Autres Expériences}

\cventry{oct 2016 -- jan 2021}{Lead engineer}{ARB-C (Association des Résidents Bétahon-Cromenach)}{}{France}{
Étude, conception, implémentation, déploiement et maintenance d'un site web pour une association locale d'habitants en Bretagne.
\begin{itemize}
\setlength{\itemindent}{0.25in}
\item Implémentation d'un blog utilisant le framework Seaside.
\item Hébergement, déploiement et maintenance de la solution via Docker.
\end{itemize}
}

\cventry{sept 2013 -- juin 2014}{Responsable de projet IT}{Ironmines (Coupe de France de Robotique)}{Douai}{France}{
Équipe de l'École des Mines Douai pour la Coupe de France de Robotique 2014.
\begin{itemize}
\setlength{\itemindent}{0.25in}
\item Conception et implémentation d'une plateforme robotique autonome utilisant le middleware ROS.
\item Implémentation de plans et routines pour relever les différents défis de la compétition.
\end{itemize}
}

\section{Langues}
\cvdoubleitem{Français}{Langue maternelle}{Anglais}{Courant}

\section{Compétences Informatiques}

\cvcomputer{Langages OO}{Smalltalk, Java, PHP, C++, CLOS, Python, Ruby}{Autres langages}{C, Bash, Perl, Lisp, Assembleur ARM}
\cvcomputer{Réingénierie}{Théorie des graphes, Visualisation scientifique, Moose}{Environnements}{Eclipse, Message-Oriented Middleware, Yii, Moose, \mbox{Jenkins}, J2EE}
\cvcomputer{Architecture}{Cloud Foundry, AWS, Docker, UML, Design Patterns, NPM, Composer, Gradle, Maven}{Gestion de projet}{Méthodologie Agile, DevOps, SAFe, Scrum, \mbox{Kanban}, TDD, ATDD, BDD}
\cvcomputer{Bases de données}{MySQL, Oracle, Elasticsearch, MongoDB}{Front-end}{HTML, CSS, JavaScript, Qt, Android, Vue.js}
\cvcomputer{Systèmes}{Windows, macOS, Linux}{Gestion de version}{Git, SVN, Mercurial}

\section{Centres d'intérêt}

\cventry{juin 2025 -- aujourd'hui}{Expédition de voile hauturière \& gestion de projet}{Projet personnel}{Irlande -- France}{}{
\begin{itemize}
    \item Planification et réalisation d'une expédition autonome autour de l'Irlande et le long de la côte Atlantique française à bord d'un monocoque en acier de 12 mètres.
    \item \textit{Compétences mises en œuvre :} Planification logistique, gestion des risques (météorologie, dépannage des systèmes), résilience et prise de décision rapide dans des environnements dynamiques.
\end{itemize}}

\cvline{Contributions Open Source} {\href{https://github.com/mattonem}{GitHub}, \href{https://www.npmjs.com/~maxmattone}{NPM}}
\cvline{Sport}{Voile (2 ans de compétition, Tall Ships Race 2014)}
\cvline{Loisirs}{Cinéma, voyages, théâtre (4 ans au conservatoire, 13 représentations en 2017), jeux de société}
\cvline{Ingénierie}{Astronomie, robotique, mécanique, voitures anciennes, électronique}

\nocite{*}
\bibliographystyle{apalike}
\bibliography{pub}



\end{document}
 
